\chapter{The Limit Order Book}\label{ch:mx:the-limit-order-book}

A trader sends an order to buy 100 shares at the best bid. The screen shows a price and a size; what it does not show is a line. In the
simulated hour this book uses throughout, 2\,834 orders joined the back of the best bid or ask: they found a median of 1\,400 shares ahead of
them, and only 13\% of them ever traded. The best price is not a place where an order sits; it is a queue that an order waits in, and the order
in which the queue formed decides who trades. This chapter describes the book as a state, as a stream of events that change the state, as a data
structure a venue and every participant must maintain, and as a queueing system whose simplest model already says how long the wait is.

\section{State: price levels and queues}

\begin{definition}[Limit order, market order]\label{def:mx:the-limit-order-book:orders}
A \emph{limit order}\index{limit order} is an order to buy at a stated price or lower, or to sell at a stated price or higher; the part that cannot
execute at once rests in the book until it executes, is cancelled or expires. A \emph{market order}\index{market order} is an order to buy or
sell a quantity at the best prices available; it never rests.
\end{definition}

\begin{definition}[Tick size, price level]\label{def:mx:the-limit-order-book:tick}
The \emph{tick size}\index{tick size} $\delta_{\mathrm{tick}}$ is the smallest price increment a venue accepts for an instrument: every limit
price is an integer multiple of it. A \emph{price level}\index{price level} is the set of resting orders on one side of the book at one price.
\end{definition}

\begin{definition}[Limit order book]\label{def:mx:the-limit-order-book:lob}
A \emph{limit order book}\index{limit order book} is the set of resting \omterm{def:mx:the-limit-order-book:orders}{limit orders} of one instrument on one venue, organised by side and by
\omterm{def:mx:the-limit-order-book:tick}{price level}; within a level, by priority. Its best bid $b_t$ and best ask $a_t$ are the highest resting buy price and the lowest resting sell price;
$Q^b_t$ and $Q^a_t$ are the displayed quantities at them.
\end{definition}

A book is \emph{crossed} if $b_t\ge a_t$; a continuously trading book never is, because an order that would cross executes instead. The
spread $s_t = a_t - b_t$ is at least one tick. Three views of the same state appear in every later chapter (One Quant Book 1, chapter 28): level 1 is
$(b_t, Q^b_t, a_t, Q^a_t)$; level 2 adds the aggregated size at each deeper price; level 3 is the book itself, every order with its reference, in
queue order (\cref{fig:mx:the-limit-order-book:book}).

\begin{omfigure}
\begin{tikzpicture}[font=\small,x=1cm,y=0.62cm]
\draw[->,thick] (0,0) -- (0,7.4) node[above]{price};
\foreach \p/\lab in {1/99.97,2/99.98,3/99.99,4/100.00,5/100.01,6/100.02,7/100.03} {\draw (-0.08,\p) -- (0.08,\p); \node[left,font=\footnotesize] at (-0.1,\p) {\lab};}
\foreach \y/\x/\w in {4/0.4/1.2,4/1.65/0.8,4/2.5/1.6,3/0.4/2.0,3/2.45/1.0,2/0.4/1.4,2/1.85/0.6,2/2.5/1.1,1/0.4/2.4} {\draw[fill=omDef!18,draw=omDef] (\x,\y-0.32) rectangle ++(\w,0.64);}
\foreach \y/\x/\w in {5/0.4/0.9,5/1.35/1.5,6/0.4/1.8,6/2.25/0.7,7/0.4/1.3,7/1.75/1.3} {\draw[fill=omThm!18,draw=omThm] (\x,\y-0.32) rectangle ++(\w,0.64);}
\draw[fill=omExo!45,draw=omExo,thick] (4.15,3.68) rectangle ++(0.7,0.64); \node[font=\footnotesize] at (4.5,4) {new};
\node[anchor=west,font=\footnotesize,align=left] at (5.0,4) {joins the back\\of the best bid};
\draw[->,thick,omThm] (5.8,5.52) -- (0.85,5.36) node[pos=0,right,font=\footnotesize,align=left]{a market buy takes\\the best ask\\from the front};
\draw[decorate,decoration={brace,amplitude=4pt}] (-1.25,4.05) -- (-1.25,4.95) node[midway,left=5pt,font=\footnotesize]{spread};
\node[font=\footnotesize,anchor=west,omDef] at (0.4,-0.6) {bids: price priority downwards, then time};
\node[font=\footnotesize,anchor=west,omThm] at (1.2,7.75) {asks: price priority upwards, then time};
\end{tikzpicture}

\omcaption{A \omterm{def:mx:the-limit-order-book:lob}{limit order book} at level 3. Each box is a resting order, its width its size; within a \omterm{def:mx:the-limit-order-book:tick}{price level} the oldest order is at the left.
The best bid is 100.00 and the best ask 100.01, one tick apart. A new bid at 100.00 joins the back of its level; a market buy consumes the best
ask from the front.}\label{fig:mx:the-limit-order-book:book}
\end{omfigure}

\section{Events: add, cancel, modify, execute}

\begin{definition}[Order-book event]\label{def:mx:the-limit-order-book:event}
An \emph{order-book event}\index{order-book event} is one change to the book: an order added at a \omterm{def:mx:the-limit-order-book:tick}{price level}, reduced or removed by its owner
(a cancellation), replaced by a new order, or reduced or removed by an execution against an incoming order. A venue's market-by-order feed
(One Quant Book 1, chapter 19) publishes one message per event, with the order's reference.
\end{definition}

The table maps the events to the messages of Nasdaq's TotalView-ITCH 5.0 feed, whose layouts the exchange simulator
of chapter~26 keeps. Nasdaq's specification is explicit about the replace: all remaining shares of the original order are removed, and the
replacement arrives with a new order reference used from then on, which is how a feed says that the order lost its place.

\begin{center}
\small
\begin{tabular}{@{}llp{4.3cm}p{4.3cm}@{}}
\toprule
event & ITCH message & level 3 & level 2 and level 1 \\
\midrule
add & A (add order) & a new order at the back of its level & the level's size grows; the best may move \\
partial cancel & X (order cancel) & the order shrinks, keeps its place & the level's size falls \\
full cancel & D (order delete) & the order leaves & the level may vanish \\
replace & U (order replace) & old order leaves, new one at the back & two levels may change \\
execution & E / C (executed) & the front order shrinks or leaves & the best level shrinks; a trade prints \\
\bottomrule
\end{tabular}
\end{center}

Executions are few. In the simulated hour (\texttt{firm.tape}, Book~7's synthetic market, 31\,535 messages), 49.7\% of messages are adds,
47.8\% cancellations and 2.5\% executions: 40 messages for each trade. Only 3.8\% of the orders that left the book did so by trading; the median
order lived 15.5 seconds (24.9 for those that traded, 15.3 for those cancelled), and the best bid price lasted a median of 4.5 seconds before it
moved. Real books in liquid stocks run faster and cancel more (chapter~3 sets the simulated facts beside published ones); what matters here is the
proportion: the book is mostly orders that will never trade.

\section{Priority and what resets it}

A \omterm{def:mx:the-limit-order-book:engine}{matching engine} (\cref{def:mx:the-limit-order-book:engine}) ranks resting orders by price, then by a rule within the level; price-time priority,
first come first served, is the common rule for equities (One Quant Book 1, chapter 19, gives the pro-rata and top-order alternatives of futures
venues). Priority is a property of an order's \emph{arrival}, and venues must say which changes count as a new arrival. The rules are close
across venues:
\begin{itemize}
\item \textbf{A reduction keeps the place.} A partial cancellation (ITCH X) leaves the order where it was: it takes away only its own shares.
\item \textbf{An increase or a price change loses it.} CME Globex lists the modifications that change an order's priority: an increase of quantity,
the price, the order type, the stop trigger price, the minimum fill quantity and the account. On Nasdaq a cancel-replace arrives on the feed as
an Order Replace with a new reference.
\item \textbf{A refreshed iceberg slice goes to the back} (chapter~2): showing more of a hidden reserve is an increase.
\end{itemize}
The reason is fairness to the orders behind: if a size increase kept its place, a trader could stand at the front with one lot and grow it
whenever a large order arrived. The consequence for a trader is that a queue position is an asset (chapter~6 prices it) that a careless
modification destroys.

\begin{definition}[Matching engine]\label{def:mx:the-limit-order-book:engine}
A \emph{matching engine}\index{matching engine} is the venue's program that holds the \omterm{def:mx:the-limit-order-book:lob}{limit order books}, applies each incoming message in the
order it arrives, executes incoming orders against resting ones by the priority rule, and publishes the resulting events and execution reports.
\end{definition}

\section{The book as a data structure}

A book has three operations on its hot path: find an order by reference (every cancel and execution names one), find the best price (every
aggressive order starts there), and walk a level in priority order (every execution). \texttt{firm.lob}, the running project's reference book,
keeps a hash map from reference to order, a map from price to level on each side, and two first-in-first-out queues per level, displayed orders
first (\cref{fig:mx:the-limit-order-book:ds}).

\begin{omfigure}
\begin{tikzpicture}[font=\footnotesize,box/.style={draw,rounded corners=2pt,minimum height=0.55cm,inner sep=3pt}]
\node[box,fill=omDef!12,minimum width=2.5cm,align=center] (h) at (0,2.6) {hash map\\ reference $\to$ order};
\node[box,fill=omDef!12,minimum width=2.5cm,align=center] (m) at (0,0) {bids: price $\to$ level\\ (sorted)};
\foreach \i/\p in {0/100.00,1/99.99,2/99.98} {\node[box,minimum width=1.3cm] (p\i) at (3.3,0.9-0.9*\i) {\p};}
\foreach \k in {0,1,2} {\node[box,fill=omDef!18] (o\k) at (5.2+1.4*\k,0.9) {\#\the\numexpr 11+\k\relax};}
\draw[->] (p0) -- (o0); \draw[->] (o0) -- (o1); \draw[->] (o1) -- (o2);
\node[box,fill=gray!15] (hid) at (9.6,0.9) {hidden};
\draw[->] (o2) -- (hid);
\draw[->] (m.east) -- (p0.west); \draw[->] (m.east) -- (p1.west); \draw[->] (m.east) -- (p2.west);
\draw[->,dashed,omDef] (h.east) to[out=0,in=100] (o1.north);
\node[anchor=west,align=left] at (5.0,-0.6) {level 100.00: displayed queue in time order,\\ then the hidden queue};
\end{tikzpicture}

\omcaption{The structures of \texttt{firm.lob}. A cancellation finds its order through the hash map (dashed); an aggressive order finds the best
level through the sorted map and walks its queues from the front.}\label{fig:mx:the-limit-order-book:ds}
\end{omfigure}

\omcode{code/firm/lob/firm_lob.py}{60}{91}{The core of \texttt{firm.lob}: add, find, remove and reduce.}\label{lst:mx:the-limit-order-book:lob}

\begin{proposition}[Costs of the book operations]\label{prop:mx:the-limit-order-book:cost}
With $L$ non-empty levels on a side and $m$ orders at a level, the structure of \cref{fig:mx:the-limit-order-book:ds} finds an order in $O(1)$ expected
time, adds an order in $O(\log L)$ when it creates a level and $O(1)$ otherwise, finds the best price in $O(1)$ (the end of the sorted map), and
removes an order in $O(\log L+m)$ worst case. A price-indexed array covering $W$ ticks, with a pointer to the best index, adds and removes in $O(1)$
and moves the best pointer in $O(g)$, where $g$ is the number of empty ticks it crosses.
\end{proposition}

\begin{proof}
Hash lookup is expected constant. A sorted map (a balanced tree, or a sorted list with binary search as in the Python reference) inserts a new key
in $O(\log L)$. Removal locates the order in its queue, linear in $m$ unless each order keeps a link to its neighbours (an intrusive doubly linked list,
the production choice, makes it $O(1)$). In the array, a level is an index computed from the price; after a removal empties the best level, the
pointer moves to the next non-empty index, $g$ steps.
\end{proof}

Production books for liquid instruments use the array: prices cluster within a few hundred ticks of the best, $g$ is small, and every
operation touches memory at a computed address (One Quant Book 13 builds one). \texttt{firm.lob} carries both: its \texttt{ArrayBook} and its
tree-keyed \texttt{MessageBook} rebuilt the simulated hour's book message by message and agreed on all 3\,154 level-2 snapshots compared. Its C++20
and Rust twins rebuild a shared five-minute fixture and reproduce every one of its 147 snapshots.

\section{The book as a queueing system}

Seen from one order, the best level is a queue in which it waits: orders ahead leave by execution (from the front) or by cancellation (from
anywhere), and the order itself may be cancelled before its turn comes. The simplest model of the whole book treats every flow as random.

\begin{definition}[Zero-intelligence model]\label{def:mx:the-limit-order-book:zi}
A \emph{zero-intelligence model}\index{zero-intelligence model} of a \omterm{def:mx:the-limit-order-book:lob}{limit order book} lets \omterm{def:mx:the-limit-order-book:orders}{limit orders}, \omterm{def:mx:the-limit-order-book:orders}{market orders} and cancellations arrive as
independent Poisson processes whose rates may depend on the distance from the best prices, with no trader choosing anything; its spread,
depth, volatility and fill probabilities follow from the rates alone.
\end{definition}

Smith, Farmer, Gillemot and Krishnamurthy (2003) built the model with \omterm{def:mx:the-limit-order-book:orders}{limit orders} placed at random over a band of prices, \omterm{def:mx:the-limit-order-book:orders}{market orders} and
cancellations, and derived by dimensional analysis how the spread, depth, volatility and impact scale with the rates; with all the rates measurable
from data, it has no free parameters. Cont, Stoikov and Talreja (2010) wrote the book as a set of birth--death queues (One Quant Book 4, chapter~8) with
cancellations at a rate proportional to the queue's size, fitted it to Tokyo Stock Exchange data and computed, among other conditional
probabilities, that of executing an order at the bid before the ask moves. The following result is the building block.

\begin{proposition}[Fill probability behind $n$ orders]\label{prop:mx:the-limit-order-book:fill}
An order joins a \omterm{def:mx:the-limit-order-book:tick}{price level} behind $n$ orders of one unit each. \omterm{def:mx:the-limit-order-book:orders}{Market orders} arrive at rate $\mu$ and each takes one unit from the front; each
order ahead is cancelled at rate $\theta$; the order itself is cancelled at rate $\nu$. All clocks are independent and exponential. The probability
that it fills before it is cancelled is
\[
P_n=\frac{\mu}{\mu+\nu}\prod_{k=1}^{n}\frac{\mu+k\theta}{\mu+k\theta+\nu},
\]
and, when $\nu=0$, the expected time until it fills is $1/\mu+\sum_{k=1}^{n}1/(\mu+k\theta)$.
\end{proposition}

\begin{proof}
With $k\ge1$ orders ahead, the next event is a departure from ahead (rate $\mu+k\theta$) or our cancellation (rate $\nu$); by memorylessness the
departure comes first with probability $(\mu+k\theta)/(\mu+k\theta+\nu)$, and the queue then restarts from $k-1$. At the front ($k=0$) the next
\omterm{def:mx:the-limit-order-book:orders}{market order} fills us with probability $\mu/(\mu+\nu)$. The stages are independent, so the probabilities multiply. With $\nu=0$ the stage
durations are exponential with means $1/(\mu+k\theta)$ and $1/\mu$; the expected total is their sum.
\end{proof}

\omcode{code/microstructure/01-the-limit-order-book/python/mx_lob.py}{115}{121}{The proposition in code.}\label{lst:mx:the-limit-order-book:fill}

The proposition holds in a queue where nobody jumps ahead. In a real book, and in the \omterm{def:mx:the-limit-order-book:zi}{zero-intelligence model}, someone does: a buyer who places
a bid one tick above ours becomes the best bid, and the market sells go to it. \Cref{fig:mx:the-limit-order-book:fill} measures the gap. The chapter's
zero-intelligence market (unit orders, \omterm{def:mx:the-limit-order-book:orders}{limit orders} at rate 0.5 per price per side over 20 prices, \omterm{def:mx:the-limit-order-book:orders}{market orders} at rate 2 per side, cancellations at
0.05 per order) tracks every bid that joins the back of the best bid. With price improvement switched off, its fill rates match the proposition with
$\mu=2$, $\theta=\nu=0.05$: behind five orders 0.871 measured against 0.870. With improvement on, as in the original model, they fall to 0.672 behind
five and 0.535 behind eight: between a quarter and a third of the fills the proposition promises are taken by orders that arrive
later at a better price. Switching
improvement off has a price of its own: the spread, 2.55 ticks with it, grows without bound without it, because nothing ever narrows it. Price
improvement is what keeps the spread finite and what makes queue position less than a guarantee.

\begin{omfigure}
\begin{tikzpicture}
\begin{axis}[name=fill,width=7.2cm,height=5.4cm,title={\small fill probability},xlabel={\small orders ahead $n$},
  tick label style={font=\footnotesize},xtick={1,2,3,4,5,6,7,8},ymin=0.4,ymax=1.0,grid=major,grid style={gray!25},
  legend columns=1,legend style={font=\footnotesize,at={(0.5,-0.3)},anchor=north,draw=none,fill=none},table/col sep=comma]
\addplot[black,thick] table[x=n,y=formula]{figdata/microstructure/01-the-limit-order-book/fill.csv};
\addplot[omDef,thick,mark=o,only marks] table[x=n,y=zi_noimprove]{figdata/microstructure/01-the-limit-order-book/fill.csv};
\addplot[omThm,thick,mark=square*] table[x=n,y=zi]{figdata/microstructure/01-the-limit-order-book/fill.csv};
\legend{proposition ($\mu=2$; $\theta=\nu=0.05$),zero intelligence without improvement,zero intelligence with improvement}
\end{axis}
\begin{axis}[at={(fill.east)},anchor=west,xshift=1.3cm,width=7.2cm,height=5.4cm,title={\small mean displayed depth (shares)},
  xlabel={\small ticks from the best},tick label style={font=\footnotesize},xtick={0,1,2,3,4,5,6,7,8,9},ymin=0,ymax=1700,grid=major,
  grid style={gray!25},ybar,bar width=4pt,legend columns=2,legend style={font=\footnotesize,at={(0.5,-0.3)},anchor=north,draw=none,fill=none},
  table/col sep=comma]
\addplot[fill=omDef!60,draw=omDef] table[x=k,y=bid]{figdata/microstructure/01-the-limit-order-book/depth.csv};
\addplot[fill=omThm!60,draw=omThm] table[x=k,y=ask]{figdata/microstructure/01-the-limit-order-book/depth.csv};
\legend{bids,asks}
\end{axis}
\end{tikzpicture}

\omcaption{Left: the probability that an order at the back of the best bid fills before it is cancelled, against the orders ahead of it: the
proposition, and the zero-intelligence market with and without price improvement (400\,000 events each; beyond eight orders ahead
the tracked orders are too few, under a hundred per point). Right: the mean displayed depth at each
distance from the best price in the simulated hour of \texttt{firm.tape}, sampled every second; it peaks one tick behind the best. Data:
\texttt{mx\_lob.zi\_market}, \texttt{mx\_lob.depth\_profile}.}\label{fig:mx:the-limit-order-book:fill}
\end{omfigure}

The right panel of \cref{fig:mx:the-limit-order-book:fill} shows a shape chapter~3 will find in real books: displayed depth is not largest at the best
price but one tick behind it, because the best level is where executions and the most cancellations happen.

\section{Tutorial: rebuild a book and time its queue}\label{tut:mx:the-limit-order-book:tut}

\begin{tutorial}
\textbf{Goal.} Rebuild the simulated hour's book in two structures, measure its event mix and its queues, and test the fill-probability
proposition in a zero-intelligence market. \textbf{End state:} the numbers of section~2, the two structures agreeing, and
\cref{fig:mx:the-limit-order-book:fill}.
\begin{tutsteps}
\item \textbf{Rebuild.} \texttt{mx\_lob.tree\_vs\_array(session())} replays 31\,535 messages into a \texttt{MessageBook} and an \texttt{ArrayBook} and
compares their ten best levels every tenth message: 3\,154 comparisons, no mismatch.
\item \textbf{Measure.} \texttt{event\_mix}, \texttt{best\_survival} and \texttt{joiners} give the shares of adds, cancels and executions, the lives of
orders, how long a best price lasts and what becomes of orders that join the best queue.
\item \textbf{Queue.} Evaluate \texttt{fill\_probability} and check it against \texttt{queue\_mc}; run \texttt{zi\_market(improve=False)} and
\texttt{zi\_market(improve=True)}; draw the figure with \texttt{fig\_lob.py}.
\item \textbf{Twins.} Run \texttt{make\_lob\_fixture.py}; build \texttt{cpp/firm\_lob\_test.cpp} and \texttt{cargo test} in \texttt{rust/}: both reproduce
the Python book's 147 snapshots.
\end{tutsteps}
\textbf{What to change next.} Give the zero-intelligence market's orders sizes drawn from a geometric law and see how far the proposition, counted in
orders, drifts from the fill rate; halve the cancellation rate and watch the spread and the depth at the best.
\end{tutorial}

\section{Build: the reference order book}\label{bld:mx:the-limit-order-book:lob}

\begin{build}
\bfield{Purpose} The book every later component of the firm stores orders in: the exchange simulator's engine (chapter~26), the participants'
own views of the market, and the tests of Book~13's production books.
\bfield{Interface} \texttt{Order(ref, side, price, qty, visible, seq, owner)}; \texttt{LimitOrderBook}: \texttt{add}, \texttt{remove}, \texttt{reduce},
\texttt{get}, \texttt{best}, \texttt{best\_visible}, \texttt{prices}, \texttt{level\_orders}, \texttt{top}, \texttt{depth}, \texttt{check};
\texttt{MessageBook.apply(kind, ref, side, price, qty, new\_ref)} for A, X, D, E, C and U; \texttt{ArrayBook(lo, hi, tick)}; C++20
\texttt{firm::lob::MessageBook} and Rust \texttt{firm\_lob::MessageBook} with the same \texttt{l2\_lines}.
\bfield{Rules} Displayed before hidden, then time, at each price; a reduction keeps the place; a replace arrives under a new reference at the back;
the book stores and never matches; \texttt{check()} verifies that every order sits in exactly its level.
\bfield{Acceptance tests} \texttt{code/firm/lob/tests/}: priority by hand; the two Python structures agreeing on a simulated session; the replace;
the fixture's 147 level-2 snapshots in C++20 and Rust.
\bfield{Stretch} An intrusive linked list per level for $O(1)$ removal; an array book that recentres when the price leaves its band.
\end{build}

\begin{omsources}
\item M.~D.~Gould, M.~A.~Porter, S.~Williams, M.~McDonald, D.~J.~Fenn and S.~D.~Howison, ``Limit order books'', \emph{Quantitative Finance}
13(11), 2013.
\item E.~Smith, J.~D.~Farmer, L.~Gillemot and S.~Krishnamurthy, ``Statistical theory of the continuous double auction'', \emph{Quantitative
Finance} 3(6), 2003.
\item R.~Cont, S.~Stoikov and R.~Talreja, ``A stochastic model for order book dynamics'', \emph{Operations Research} 58(3), 2010.
\item Nasdaq, \emph{TotalView-ITCH 5.0} specification; CME Group, client systems documentation, ``Order functionalities''.
\end{omsources}

\section{Exercises}

\begin{exercise}[$\star$]\label{exo:mx:the-limit-order-book:1}
The best bid is 100.00 with 700 shares in three orders of 300, 200 and 200; the best ask is 100.01. A market sell of 450 shares arrives. What trades,
and what is the new level 1?
\end{exercise}

\begin{exercise}[$\star$]\label{exo:mx:the-limit-order-book:2}
An order joins behind three unit orders; \omterm{def:mx:the-limit-order-book:orders}{market orders} arrive at rate 1 per second, each order ahead cancels at 0.1 per second, and the order's
owner cancels it at 0.2 per second. What is its fill probability?
\end{exercise}

\begin{exercise}[$\star$]\label{exo:mx:the-limit-order-book:3}
A trader first reduces an order from 500 to 300 shares, then raises it back to 500. What happens to its queue position at each step on CME Globex,
and how does each step appear in an ITCH feed?
\end{exercise}

\begin{exercise}[$\star\star$]\label{exo:mx:the-limit-order-book:4}
A price-indexed array covers the prices within 1\% of 100.00 for a stock with a one-cent tick. How many slots does it need, and what happens on a
day the stock moves 3\%?
\end{exercise}

\begin{exercise}[$\star\star$]\label{exo:mx:the-limit-order-book:5}
With no cancellation of its own, an order waits behind ten unit orders; \omterm{def:mx:the-limit-order-book:orders}{market orders} arrive at 0.5 per second and each order ahead cancels at 0.02 per
second. What is its expected wait?
\end{exercise}

\begin{exercise}[$\star\star$]\label{exo:mx:the-limit-order-book:6}
Why is displayed depth in \cref{fig:mx:the-limit-order-book:fill} larger one tick behind the best than at the best?
\end{exercise}

\begin{exercise}[$\star\star\star$]\label{exo:mx:the-limit-order-book:7}
\emph{Coding.} Rerun \texttt{zi\_market} with \omterm{def:mx:the-limit-order-book:orders}{market orders} at rate 1 instead of 2. What happens to the spread, and to the gap between the
proposition and the measured fill probability behind five orders? Explain both.
\end{exercise}

\begin{exercise}[$\star\star\star$]\label{exo:mx:the-limit-order-book:8}
\emph{Find the flaw.} ``Our order has been at the front of the best bid for ten seconds, so by the proposition it will fill with probability
$\mu/(\mu+\nu)$ whatever happens next.''
\end{exercise}

\section{Problem: The Line at the Best Bid}

\begin{problem}[{Weekend problem --- the line at the best bid}]\label{pb:mx:the-limit-order-book:1}
A desk asks how long its passive orders wait and how many of them trade. You have the simulated hour of \texttt{firm.tape} and the chapter's
zero-intelligence market.

\medskip\noindent\textbf{Part I --- The stream.}
\begin{enumerate}
\item How many messages does the hour hold, and in what proportions of adds, cancels and executions?
\item How many messages are there per trade, and what share of orders leave the book by trading?
\item How long does the median order live, and does an order that trades live longer than one that is cancelled?
\item How long does a best bid price last, in the median?
\end{enumerate}

\medskip\noindent\textbf{Part II --- The queue at the best.}
\begin{enumerate}[resume]
\item How many orders join the back of the best bid or ask, and how many shares do they find ahead?
\item What share of them ever trade, and what share fill in full?
\item Where does mean displayed depth peak, and why there?
\item Which of these numbers would change if cancellations were slower?
\end{enumerate}

\medskip\noindent\textbf{Part III --- The model.}
\begin{enumerate}[resume]
\item State the fill-probability proposition and its assumptions.
\item Compute $P_1$, $P_5$ and $P_{12}$ for $\mu=2$, $\theta=\nu=0.05$.
\item Check $P_5$ by simulation of the stylised queue.
\item What is the expected wait behind five orders with $\nu=0$?
\end{enumerate}

\medskip\noindent\textbf{Part IV --- The market.}
\begin{enumerate}[resume]
\item What fill rates does the zero-intelligence market give without price improvement, and why do they match the proposition?
\item What does it give with improvement, behind one, five and eight orders?
\item State the \emph{named result}: the fill probability of an order at the back of the best bid against the orders ahead, and the share of it
that price improvement takes away.
\item Why does the market without improvement have an unbounded spread?
\item What is the zero-intelligence market's spread and mean best-level depth with improvement?
\item Which assumption of the proposition fails most in a real book, and in which direction does it bias the answer?
\item What would you measure in real data to estimate $\mu$, $\theta$ and $\nu$?
\item In one sentence: what is a queue position worth, before chapter~6 prices it?
\end{enumerate}
\end{problem}

\section{Interview questions}

\begin{interviewq}[$\star$ \iqroles{developer}]\label{iq:mx:the-limit-order-book:1}
Design the data structures of a \omterm{def:mx:the-limit-order-book:lob}{limit order book}. What are the costs of add, cancel, execute and best-price lookup?
\end{interviewq}

\begin{interviewq}[$\star$ \iqroles{trader, researcher}]\label{iq:mx:the-limit-order-book:2}
Why does increasing an order's size lose its queue position, while decreasing it does not?
\end{interviewq}

\begin{interviewq}[$\star\star$ \iqroles{researcher}]\label{iq:mx:the-limit-order-book:3}
You are fifth in a queue. \omterm{def:mx:the-limit-order-book:orders}{Market orders} arrive at rate $\mu$, orders ahead cancel at rate $\theta$ each. What is the probability you fill before you
cancel at rate $\nu$?
\end{interviewq}

\begin{interviewq}[$\star\star$ \iqroles{developer}]\label{iq:mx:the-limit-order-book:4}
A price-indexed array book is fast. When does it fail, and what do you do about it?
\end{interviewq}

\begin{interviewq}[$\star\star$ \iqroles{researcher, trader}]\label{iq:mx:the-limit-order-book:5}
Most orders in a modern book are cancelled. Is that a problem for the market?
\end{interviewq}

\begin{interviewq}[$\star\star\star$ \iqroles{researcher}]\label{iq:mx:the-limit-order-book:6}
What does a \omterm{def:mx:the-limit-order-book:zi}{zero-intelligence model} of the order book get right, and what does it get wrong?
\end{interviewq}
